\subsection{The evaluation relation and Reed--Solomon proximity}\label{sec:security:relations}

\begin{definition}[Kernel encoding as a word]\label{def:security:encoding}
Retain the finite field $\F$, integer $m\geq1$, $N=2^m$, and evaluation domain $\mathcal D\subseteq\F$ with $M:=|\mathcal D|\geq N$. For each table $f:I_m\to\F$ define
\begin{equation}\label{eq:security:encoding}
 \Enc_\mathcal D(f):\mathcal D\to\F,\qquad
 \Enc_\mathcal D(f)(\xi):=U_f^{(m)}(\xi).
\end{equation}
The corresponding Reed--Solomon code is
\[
 \mathcal C_{\mathcal D,N}
 :=\{\Enc_\mathcal D(f):f\in\F^{I_m}\}
 \subseteq\F^\mathcal D.
\]
\end{definition}

\begin{lemma}[Injectivity of kernel encoding]\label{lem:security:encoding-injective}
If $\Enc_\mathcal D(f)=\Enc_\mathcal D(g)$ as words, then $f=g$.
\end{lemma}
\begin{proof}
Let $P:=U_f^{(m)}-U_g^{(m)}$. By Lemma~\ref{lem:intro:degree} and the linear definition of $U_f$, $P$ has degree less than $N$, or is zero. The equality of words means that $P$ vanishes at every $\xi\in\mathcal D$. If $P$ were nonzero, Lemma~\ref{lem:history:rootcount} would imply at most $N-1$ roots, contradicting $|\mathcal D|=M\geq N$. Thus $P=0$. By the kernel basis theorem, Proposition~\ref{prop:intro:basis}, its coordinates satisfy $f(b)-g(b)=0$ for every $b$, hence $f=g$.
\end{proof}

\begin{definition}[Hamming distance]\label{def:security:hamming}
For two words $w,w':\mathcal D\to\F$ define their Hamming distance by
\[
 d_H(w,w'):=\bigl|\{\xi\in\mathcal D:w(\xi)\neq w'(\xi)\}\bigr|
       \in\{0,\ldots,M\}.
\]
The relative Hamming distance is $d_H(w,w')/M$.
\end{definition}

\begin{lemma}[Triangle inequality for Hamming distance]\label{lem:security:hamming-triangle}
For every three words $w_1,w_2,w_3:\mathcal D\to\F$,
\[
 d_H(w_1,w_3)\leq d_H(w_1,w_2)+d_H(w_2,w_3).
\]
\end{lemma}
\begin{proof}
If $w_1(\xi)\neq w_3(\xi)$, then at least one of
$w_1(\xi)\neq w_2(\xi)$ and $w_2(\xi)\neq w_3(\xi)$ holds; otherwise transitivity of equality would give $w_1(\xi)=w_3(\xi)$. Thus the set of coordinates where $w_1$ and $w_3$ differ is contained in the union of the two sets where consecutive words differ. Taking cardinalities and using the union bound for finite sets proves the inequality.
\end{proof}

\begin{theorem}[Reed--Solomon minimum distance]\label{thm:security:rs-distance}
For any distinct tables $f,g:I_m\to\F$,
\begin{equation}\label{eq:security:rs-distance}
 d_H(\Enc_\mathcal D(f),\Enc_\mathcal D(g))\geq M-N+1.
\end{equation}
\end{theorem}
\begin{proof}
By Proposition~\ref{prop:intro:basis}, $f\neq g$ implies
$P:=U_f^{(m)}-U_g^{(m)}\neq0$. This polynomial has degree at most $N-1$. The two encoded words agree precisely at positions $\xi\in\mathcal D$ with $P(\xi)=0$. Lemma~\ref{lem:history:rootcount} shows that there are at most $N-1$ such positions. Consequently there are at least $M-(N-1)=M-N+1$ differing positions, which is~\eqref{eq:security:rs-distance}.
\end{proof}

\begin{definition}[Exact evaluation relation]\label{def:security:exactrel}
An \emph{exact oracle evaluation instance} is $(w;z,v)$ with
$w:\mathcal D\to\F$, $z\in\F^m$, and $v\in\F$. Its witness relation is
\begin{equation}\label{eq:security:exact-rel}
 \mathcal R_{\rm eval}:=
 \left\{((w;z,v),f):
     w=\Enc_\mathcal D(f)\ \land\ v=\widetilde f(z)\right\}.
\end{equation}
The semicolon separates the oracle part $w$ from the explicit part $(z,v)$; the verifier may read $w$ only through its oracle.
\end{definition}

\begin{definition}[Proximity evaluation relation]\label{def:security:proxrel}
For an integer $\tau$ satisfying $0\leq\tau\leq M$, define
\begin{equation}\label{eq:security:prox-rel}
 \mathcal R_{\tau}:=
 \left\{((w;z,v),f):
       d_H(w,\Enc_\mathcal D(f))\leq\tau\
       \land\ v=\widetilde f(z)\right\}.
\end{equation}
Its language is
\[
 \mathcal L_\tau:=\{(w;z,v):\exists f,\ ((w;z,v),f)\in\mathcal R_\tau\}.
\]
The relation $\mathcal R_0$ equals $\mathcal R_{\rm eval}$.
\end{definition}

\begin{lemma}[Uniqueness of a nearby codeword]\label{lem:security:unique-near}
Assume
\[
 2\tau<M-N+1.
\]
For every received word $w:\mathcal D\to\F$ there exists \emph{at most one} table
$f:I_m\to\F$ such that $d_H(w,\Enc_\mathcal D(f))\leq\tau$.
\end{lemma}
\begin{proof}
Suppose $f$ and $g$ both satisfy the distance condition. By Lemma~\ref{lem:security:hamming-triangle},
\[
 d_H(\Enc_\mathcal D(f),\Enc_\mathcal D(g))
 \leq d_H(\Enc_\mathcal D(f),w)+d_H(w,\Enc_\mathcal D(g))
 \leq2\tau.
\]
If $f\neq g$, Theorem~\ref{thm:security:rs-distance} gives a contradictory lower bound $M-N+1>2\tau$. Therefore $f=g$.
\end{proof}

\begin{corollary}[At most one nearby evaluation value]\label{cor:security:unique-value}
Under the inequality $2\tau<M-N+1$, for every pair $(w,z)$ there exists at most one $v\in\F$ such that $(w;z,v)\in\mathcal L_\tau$.
\end{corollary}
\begin{proof}
Suppose $v$ and $v'$ both occur. By the definition of $\mathcal L_\tau$, there are corresponding tables $f,g$ within distance $\tau$ of $w$. By Lemma~\ref{lem:security:unique-near}, $f=g$. The relation requires $v=\widetilde f(z)$ and $v'=\widetilde g(z)$, so $v=v'$.
\end{proof}

\begin{definition}[Oracle evaluation soundness]\label{def:security:eval-sound}
For the fixed proximity threshold $\tau$ and an oracle IOP $\Pi=(P,V)$, define
\[
 \mathsf{Sound}_\tau(\Pi)
 :=\sup_{\substack{(w;z,v)\notin\mathcal L_\tau\\P^*\ {\rm causal}}}
  \Pr\bigl[V^{P^*}(w;z,v)\text{ accepts}\bigr],
\]
where the prover's randomness and the honest verifier's challenges are independent, and the supremum over an empty class is defined to be $0$. The statement $\mathsf{Sound}_\tau(\Pi)\leq\varepsilon$ is \emph{$\varepsilon$-evaluation soundness at proximity $\tau$}.
\end{definition}

\begin{remark}[What the Reed--Solomon lemmas do and do not prove]
Theorem~\ref{thm:security:rs-distance} and Lemma~\ref{lem:security:unique-near} establish only code-theoretic uniqueness. They do not prove the soundness of KBFold itself: that requires bounds on the probability of passing its fold and oracle queries. The parent KBFold paper supplies such a soundness analysis in the unique-decoding regime~\cite{MkhidaKBFold2026}, which must be revisited when the polynomial degree increases under zero-knowledge masking.
\end{remark}
